\chapter{Dane i opis zjawiska}

\section{Charakterystyka środowiska League of Legends}

\textit{League of Legends} jest drużynową grą strategiczną typu MOBA, w której dwa pięcioosobowe
zespoły rywalizują na symetrycznej mapie. Celem gry jest zniszczenie głównej struktury przeciwnika,
ale droga do tego celu obejmuje wiele mierzalnych elementów: kontrolę zasobów, przewagę ekonomiczną,
walki drużynowe, niszczenie wież, zdobywanie smoków i Baronów oraz wykorzystanie indywidualnych
umiejętności zawodników. Z perspektywy tej pracy najważniejsze jest to, że część tych informacji
można obserwować historycznie i przekształcać w cechy opisujące styl oraz formę drużyny przed
kolejnym meczem.

Profesjonalny mecz może być rozgrywany jako pojedyncza gra Bo1 albo jako seria Bo3 lub Bo5. Format
serii ma znaczenie probabilistyczne: w dłuższej serii przypadkowy wynik jednej mapy ma mniejszy
wpływ na końcowe rozstrzygnięcie niż w Bo1. Z tego powodu format meczu nie jest w pracy traktowany
wyłącznie jako opis organizacyjny, lecz jako informacja strukturalna wykorzystywana później w modelu
W20-Binomial.

Rozgrywki odbywają się w wielu regionach i na różnych poziomach sportowych. Obok lig najwyższego
szczebla, takich jak LCK, LPL, LEC i LCS, występują ligi regionalne, akademickie oraz turnieje
międzynarodowe. Taki zbiór danych daje nam dużą wariancję poziomów rozgrywek, oraz bogatszy zbiór do
uczenia modelu.

Istotną cechą środowiska jest również zmienność w czasie. Aktualizacje gry wpływają na siłę postaci,
tempo rozgrywki i opłacalność strategii, a transfery zawodników zmieniają realną siłę drużyn.
Konsekwencją dla modelowania jest konieczność chronologicznego przetwarzania danych, oraz szybkie
adaptoanie się do zmian sił drużyn.

\section{Dlaczego profesjonalne League of Legends?}

Wybór profesjonalnych meczów \textit{League of Legends} jako obszaru badawczego wynika nie tylko z
popularności gry, lecz przede wszystkim z jej przydatności jako przypadku testowego dla modeli
przedmeczowych. Dane łączą trzy elementy rzadko występujące jednocześnie: długą historię wyników,
identyfikowalne składy zawodników oraz zewnętrzny benchmark w postaci kursów bukmacherskich. Dzięki
temu można badać zarówno modele rankingowe, jak i jakość końcowych prawdopodobieństw względem rynku.

Z punktu widzenia pracy inżynierskiej jest to środowisko pośrednie między klasycznym sportem
drużynowym a systemem cyfrowym. Podobnie jak w tradycyjnych ligach występują tu forma, siła
zawodników, poziom przeciwnika i format rozgrywek. Jednocześnie gra podlega regularnym
aktualizacjom, a składy mogą zmieniać się szybciej niż w wielu sportach tradycyjnych. Taka
kombinacja uzasadnia użycie modeli aktualizowanych chronologicznie oraz ocenę ich stabilności w
czasie.

Istotny jest także charakter istniejących badań. Część prac dotyczących \textit{League of Legends}
koncentruje się na predykcji w trakcie trwania gry, gdy znane są już statystyki meczowe, takie jak
złoto, zabójstwa, obrażenia lub inne zmienne opisujące aktualny stan rozgrywki. Przykładowo Lin
pokazał, że predykcja oparta na informacjach przedmeczowych była znacznie słabsza niż predykcja
wykorzystująca statystyki z trwającej gry \parencite{lin2016lol}. Podobnie praca Juniora i Campelo
dotyczy predykcji w czasie rzeczywistym i oceny modeli na różnych etapach meczu, osiągając najwyższe
wyniki w środkowych fazach rozgrywki \parencite{junior2023lol}. Inne prace wykorzystują informacje
dostępne po wyborze postaci, na przykład doświadczenie zawodników na wybranych bohaterach
\parencite{do2021championExperience}.

Na tym tle niniejsza praca przyjmuje węższe i trudniejsze założenie informacyjne: model ma
przewidywać wynik profesjonalnego meczu przed jego rozpoczęciem, bez dostępu do statystyk z bieżącej
gry ani do informacji powstających po starcie spotkania. Zamiast tego wykorzystuje wyłącznie
historię wcześniejszych meczów, systemy rankingowe, miary niepewności oraz historyczne agregaty
formy. Takie ujęcie odróżnia pracę od badań typu live prediction i lepiej odpowiada zadaniu
porównania z kursem otwarcia rynku bukmacherskiego.

\section{Źródła danych}

W pracy zebrano dwa właściwe źródła danych. Pierwszym źródłem jest serwis GOL.GG (Games of League)
\parencite{golgg}, udostępniający dane sportowe dotyczące profesjonalnych meczów \textit{League of
Legends}. W dalszej części pracy stosowana jest wyłącznie skrócona nazwa GOL.GG. Dane te obejmują
między innymi datę rozegrania spotkania, nazwy drużyn, wynik meczu, format serii oraz informacje
dotyczące poszczególnych gier i zawodników.

Drugim źródłem są historyczne kursy bukmacherskie pochodzące z serwisu OddsPortal
\parencite{oddsportal}. W pracy kursy nie są wykorzystywane do budowy strategii zakładów, lecz jako
benchmark predykcyjny. Prawdopodobieństwa implikowane przez kursy bukmacherskie stanowią punkt
odniesienia dla modeli opartych wyłącznie na danych sportowych.

Pozostałe zbiory wykorzystywane w eksperymentach mają charakter pośredni. Powstały w wyniku
przetworzenia danych sportowych, na przykład przez wyznaczenie predykcji systemów rankingowych lub
obliczenie historycznych cech kroczących. Nie stanowią one dodatkowych, zewnętrznych źródeł
informacji. Rozróżnienie to jest istotne, ponieważ część badawcza pracy opiera się na dwóch źródłach
pierwotnych: GOL.GG oraz OddsPortal.

Zakres danych można podzielić na dwa poziomy. Pełny zbiór GOL.GG służy do odtwarzania historii
sportowej, generowania rankingów oraz obliczania cech historycznych. Główna ewaluacja wyników odbywa
się jednak na wspólnej próbce meczów, dla których dostępne są zarówno dane sportowe, jak i kursy
bukmacherskie. Takie ograniczenie zapewnia porównywalność modeli z rynkiem bukmacherskim.

Po przygotowaniu danych pełna historia sportowa obejmowała 40~656 meczów z lat 2013--2026 oraz
66~223 pojedyncze gry. Najliczniejszą grupę stanowiły mecze Bo1, ale w danych obecne były również
serie Bo3 i Bo5. Zbiór kursów bukmacherskich obejmował 16~375 dopasowanych meczów z okresu od 6
kwietnia 2017 roku do 22 marca 2026 roku. Z punktu widzenia końcowej ewaluacji najważniejsza jest
nie sama liczba rekordów, lecz fakt, że wspólna próbka pozwala porównać predykcje modelu sportowego
i rynku na tych samych spotkaniach.

\begin{table}[H]
    \centering
    \caption{Źródła danych i zbiory pośrednie wykorzystywane w pracy}
    \label{tab:data_sources_role}
    \begin{tabular}{lll}
        \toprule
        Element & Typ & Rola w pracy \\
        \midrule
        GOL.GG & źródło pierwotne & dane sportowe i historia meczów \\
        OddsPortal & źródło pierwotne & historyczne kursy bukmacherskie \\
        Predykcje rankingowe & zbiór pośredni & cechy wejściowe metamodelu \\
        Cechy kroczące & zbiór pośredni & historyczny kontekst formy \\
        \bottomrule
    \end{tabular}
\end{table}

\begin{table}[H]
    \centering
    \caption{Podstawowe wielkości analizowanych zbiorów danych}
    \label{tab:data_scope_summary}
    \small
    \begin{tabular}{p{5.3cm}r c p{3.2cm}}
        \toprule
        Zakres & Liczba obserwacji & Okres & Rola w analizie \\
        \midrule
        Mecze sportowe GOL.GG & 40~656 & 2013--2026 & historia sportowa i rankingi \\
        Pojedyncze gry GOL.GG & 66~223 & 2013--2026 & statystyki gry i kontekst \\
        Mecze z kursami OddsPortal & 16~375 & 2017--2026 & benchmark rynkowy \\
        Finalna próbka model--rynek & 11~609 & 2021--2026 & końcowe porównanie modeli \\
        \bottomrule
    \end{tabular}
\end{table}

\section{Zakres informacyjny danych sportowych}

Dane sportowe z GOL.GG mają strukturę hierarchiczną. Najwyższym poziomem obserwacji jest mecz, czyli
spotkanie dwóch drużyn rozstrzygane w formacie Bo1, Bo3 albo Bo5. Pojedynczy mecz może składać się z
jednej lub kilku gier. Z tego powodu w pracy rozróżniono poziom meczu, na którym definiowany jest
wynik predykcji, oraz poziom gry, na którym dostępna jest część statystyk opisujących przebieg
rozgrywki.

Na poziomie meczu dostępne są przede wszystkim: data spotkania, nazwy i identyfikatory drużyn, nazwa
turnieju, format serii, wynik końcowy oraz informacja, która drużyna wygrała. Te zmienne są
wykorzystywane do uporządkowania obserwacji w czasie, zdefiniowania zmiennej objaśnianej oraz
segmentacji wyników według formatu meczu i poziomu rozgrywek.

Na poziomie gry dostępne są informacje bardziej szczegółowe. Obejmują one między innymi składy
zawodników, wynik pojedynczej gry oraz statystyki drużynowe i zawodnicze opisujące przebieg mapy. W
tej grupie znajdują się zmienne takie jak liczba zabójstw, zadane obrażenia, udział w obrażeniach,
tempo gry oraz przewaga w złocie w określonym momencie gry. W pracy nie wykorzystano ich jako
informacji z aktualnego meczu. Zamiast tego posłużyły do budowy historycznych agregatów, obliczanych
wyłącznie na podstawie wcześniejszych występów.

\begin{table}[H]
    \centering
    \caption{Główne grupy zmiennych sportowych wykorzystywanych w pracy}
    \label{tab:sport_variable_groups}
    \small
    \begin{tabular}{p{3.8cm}p{4.6cm}p{5.0cm}}
        \toprule
        Grupa zmiennych & Przykładowe informacje & Rola w pracy \\
        \midrule
        Identyfikacja meczu & data, turniej, drużyny, identyfikator meczu & porządkowanie danych, dopasowanie do kursów, walidacja chronologiczna \\
        Wynik meczu & zwycięzca, wynik serii, format Bo1/Bo3/Bo5 & definicja zmiennej objaśnianej i segmentacja wyników \\
        Skład drużyn & zawodnicy Team 1 i Team 2, role graczy & budowa rankingów zawodniczych i porównanie z rankingami drużynowymi \\
        Statystyki gry & zabójstwa, obrażenia, tempo gry, GD@15 & źródło historycznych cech kroczących opisujących formę i styl drużyny \\
        Poziom rozgrywek & liga, turniej, segment Tier 1 lub Regional/ERL & analiza stabilności wyników w różnych segmentach sceny \\
        \bottomrule
    \end{tabular}
\end{table}

Najważniejsze jest rozróżnienie pomiędzy zmiennymi dostępnymi przed meczem a zmiennymi powstającymi
w trakcie gry. Wynik aktualnej gry, aktualne zabójstwa lub aktualna przewaga w złocie nie mogą być
użyte do predykcji przedmeczowej. Mogą natomiast zostać przekształcone w cechy historyczne, jeżeli
są obliczane wyłącznie na podstawie wcześniejszych meczów tej samej drużyny.

\section{Zakres informacyjny danych bukmacherskich}

Dane z OddsPortal opisują rynek kursów dla meczów profesjonalnych. Dla każdego dopasowanego
spotkania dostępne są kursy na zwycięstwo obu stron meczu. W pracy szczególne znaczenie mają dwie
wersje kursów: kurs otwarcia oraz kurs zamknięcia.

Kurs otwarcia reprezentuje wcześniejszą wycenę rynku i jest najbliższy momentowi, w którym model
przedmeczowy mógłby zostać użyty. Kurs zamknięcia jest ceną późniejszą, zwykle bliższą rozpoczęciu
meczu. Może więc zawierać dodatkowe informacje, korekty rynku oraz reakcję na zmiany dostępnej
wiedzy. Z tego powodu w pracy kurs zamknięcia traktowany jest jako benchmark diagnostyczny, a nie
jako informacja dostępna w tym samym momencie co pierwotna predykcja modelu.

\begin{table}[H]
    \centering
    \caption{Główne grupy zmiennych rynkowych wykorzystywanych w pracy}
    \label{tab:market_variable_groups}
    \small
    \begin{tabular}{p{3.8cm}p{4.8cm}p{4.8cm}}
        \toprule
        Grupa zmiennych & Przykładowe informacje & Rola w pracy \\
        \midrule
        Identyfikacja spotkania & data, nazwy drużyn, turniej, strony meczu & dopasowanie kursów do danych GOL.GG \\
        Kursy otwarcia & średni kurs rynkowy i kursy bukmacherów przed korektami rynku & główny benchmark przedmeczowy \\
        Kursy zamknięcia & kursy bliższe rozpoczęciu meczu & silniejszy benchmark diagnostyczny \\
        Bukmacherzy & Betclic, Betfan, eFortuna, LV BET, STS, Superbet & analiza pokrycia i marży poszczególnych źródeł kursów \\
        Prawdopodobieństwa bez marży & znormalizowane prawdopodobieństwa implikowane z kursów & porównanie z prawdopodobieństwami modeli \\
        \bottomrule
    \end{tabular}
\end{table}

\section{Definicja problemu predykcyjnego}

Problem badawczy został sformułowany jako binarna predykcja zwycięzcy meczu. Dla każdego spotkania
model przewiduje prawdopodobieństwo zwycięstwa pierwszej drużyny w uporządkowanej parze meczowej.
Zmienną objaśnianą oznaczono jako:

\begin{equation}
    y =
    \begin{cases}
        1, & \text{jeżeli pierwsza drużyna wygrała mecz,} \\
        0, & \text{jeżeli pierwsza drużyna przegrała mecz.}
    \end{cases}
\end{equation}

Modele generują wartość \(\hat{p}\), interpretowaną jako szacowane prawdopodobieństwo zdarzenia \(y
= 1\). W pracy nie ograniczono się do klasyfikacji zwycięzcy przy progu 0{,}5. Istotniejsza jest
jakość prawdopodobieństwa, ponieważ dobrze skalibrowany model powinien nie tylko wskazywać bardziej
prawdopodobną drużynę, lecz także realistycznie oceniać poziom niepewności.

\section{Kursy bukmacherskie jako benchmark probabilistyczny}

Kursy bukmacherskie zapisane w formacie dziesiętnym można przekształcić na prawdopodobieństwa
implikowane. Dla rynku dwuwynikowego, w którym dostępny jest kurs na zwycięstwo pierwszej i drugiej
drużyny, surowe prawdopodobieństwa wynikające z kursów mają postać:

\begin{equation}
    q_1 = \frac{1}{o_1}, \qquad q_2 = \frac{1}{o_2},
\end{equation}

gdzie \(o_1\) i \(o_2\) oznaczają kursy dziesiętne na odpowiednio pierwszą i drugą drużynę. Ponieważ
kursy zawierają marżę bukmachera, suma \(q_1 + q_2\) zwykle przekracza jeden. W celu uzyskania
prawdopodobieństwa porównywalnego z modelem zastosowano normalizację bez marży:

\begin{equation}
    p_1 = \frac{q_1}{q_1 + q_2}
    = \frac{1/o_1}{1/o_1 + 1/o_2}.
\end{equation}

W pracy rozróżniono dwa typy kursów. Kursy otwarcia, oznaczane jako \textit{Market Open}, traktowane
są jako operacyjny punkt odniesienia, ponieważ reprezentują informację dostępną wcześniej. Kursy
zamknięcia, oznaczane jako \textit{Market Close}, mają charakter diagnostyczny. Zawierają późniejszą
informację rynkową i dlatego nie powinny być interpretowane jako cena dostępna w momencie pierwszej
predykcji.

W zbiorze kursów dostępne były średnie kursy rynkowe oraz kursy wybranych bukmacherów. W analizie
uwzględniono przede wszystkim Betclic, Betfan, eFortuna, LV BET, STS oraz Superbet. W danych
technicznych występowały również kolumny dla Fuksiarza, jednak w analizowanej wersji zbioru nie
miały one użytecznego pokrycia i nie zostały wykorzystane w porównaniach jakości rynku. Mediana
średniego kursu otwarcia wynosiła 1{,}78 dla pierwszej strony meczu oraz 1{,}93 dla drugiej strony.
Maksymalne obserwowane wartości średnich kursów otwarcia wynosiły odpowiednio 16{,}50 oraz 19{,}00,
co pokazuje obecność wyraźnych faworytów i underdogów.

\begin{table}[H]
    \centering
    \caption{Zakres średnich kursów rynkowych w zbiorze OddsPortal}
    \label{tab:odds_value_ranges}
    \begin{tabular}{lrrrr}
        \toprule
        Zmienna & Liczba obserwacji & Minimum & Mediana & Maksimum \\
        \midrule
        Kurs otwarcia -- strona 1 & 16~316 & 1{,}00 & 1{,}78 & 16{,}50 \\
        Kurs otwarcia -- strona 2 & 16~316 & 1{,}00 & 1{,}93 & 19{,}00 \\
        Kurs zamknięcia -- strona 1 & 16~316 & 1{,}00 & 1{,}79 & 16{,}50 \\
        Kurs zamknięcia -- strona 2 & 16~316 & 1{,}01 & 1{,}93 & 16{,}06 \\
        \bottomrule
    \end{tabular}
\end{table}

\section{Przygotowanie i dopasowanie danych}

Istotnym etapem przygotowania danych było dopasowanie meczów sportowych z GOL.GG do meczów
posiadających kursy bukmacherskie z OddsPortal. Dopasowanie to wymagało uwzględnienia daty
spotkania, nazw drużyn, stron meczu oraz wyniku końcowego. Celem tego etapu było utworzenie wspólnej
próbki obserwacji, w której dla tego samego meczu dostępne są zarówno informacje sportowe, jak i
benchmark rynkowy.

Na podstawie danych sportowych wyznaczono następnie cechy używane w modelowaniu. Systemy rankingowe
były aktualizowane chronologicznie po każdym meczu, a przed kolejnym spotkaniem generowały predykcję
siły drużyn. Dodatkowo obliczono cechy historyczne opisujące ostatnią formę zespołów. Wszystkie te
zmienne powstały wyłącznie z informacji dostępnych przed analizowanym meczem.

Szczególnie ważne było również zachowanie poprawnej orientacji stron meczu. W części rekordów
kolejność drużyn w historycznym zbiorze kursów różniła się od kolejności drużyn w danych sportowych.
Z tego powodu prawdopodobieństwa rynkowe zostały wyrównane do tej samej strony, dla której modele
generowały predykcję. Rekordy, dla których nie można było jednoznacznie ustalić orientacji, zostały
wyłączone z finalnej wspólnej próbki ewaluacyjnej.

\section{Braki danych i ograniczenia próby}

Zgromadzone dane są wystarczające do budowy modeli rankingowych i probabilistycznych, ale nie
opisują pełnego kontekstu sportowego meczu. W szczególności w pracy nie wykorzystano informacji o
wyborze postaci, fazie draftu, planowanych strategiach, stanie zdrowia zawodników, zmianach
treningowych ani wewnętrznej sytuacji organizacji. Część z tych informacji może być znana
uczestnikom rynku lub analitykom drużyn, ale nie jest dostępna w uporządkowanej, historycznej
postaci dla całej analizowanej próby.

Drugim ograniczeniem jest pokrycie danych bukmacherskich. Nie wszystkie mecze z pełnej historii
GOL.GG posiadają odpowiadające im kursy w OddsPortal. W konsekwencji finalna próba ewaluacyjna jest
mniejsza od pełnej historii sportowej i obejmuje przede wszystkim mecze, dla których rynek
bukmacherski faktycznie udostępniał ceny. Oznacza to, że końcowe porównanie z rynkiem dotyczy
ważniejszej i bardziej obserwowalnej części profesjonalnej sceny, ale nie powinno być automatycznie
uogólniane na wszystkie mecze dostępne w GOL.GG.

Trzecim ograniczeniem jest nierównomierne pokrycie bukmacherów. Część firm ma dużą liczbę
obserwacji, natomiast inne występują rzadziej lub nie mają użytecznego pokrycia w analizowanej
wersji danych. Z tego powodu w głównym benchmarku wykorzystano średnie prawdopodobieństwo rynkowe po
normalizacji kursów, a nie predykcję pojedynczego bukmachera.

Ostatnim ograniczeniem jest interpretacja kursów zamknięcia. Market Close jest silnym benchmarkiem
diagnostycznym, ponieważ zawiera późniejsze informacje i reakcję rynku. Nie jest jednak równoważny
kursowi dostępnemu w momencie pierwotnej predykcji przedmeczowej. W pracy rozdzielono więc Market
Open jako wcześniejszy punkt odniesienia oraz Market Close jako benchmark ex post.

\section{Zasady unikania przecieku informacji}

Wszystkie modele oceniane w pracy wymagają zachowania porządku czasowego. System rankingowy przed
danym meczem może korzystać wyłącznie z wyników rozegranych wcześniej spotkań. Po wygenerowaniu
predykcji wynik meczu jest używany do aktualizacji ratingów, ale tylko dla kolejnych obserwacji.
Analogiczna zasada dotyczy cech historycznych, takich jak średnia forma lub statystyki kroczące.

Takie podejście odpowiada realistycznemu scenariuszowi predykcyjnemu. Model nie zna przyszłych
wyników, przyszłych składów ani późniejszych zmian kursów. W szczególności cechy opisujące przebieg
aktualnej gry, takie jak przewaga w złocie w piętnastej minucie bieżącej mapy, nie mogą być używane
bezpośrednio do predykcji przedmeczowej. Mogą natomiast zostać wykorzystane w formie historycznych
średnich kroczących, obliczonych na podstawie wcześniejszych występów drużyny.

\section{Uzasadnienie wyboru głównej próbki ewaluacyjnej}

Chociaż pełny zbiór GOL.GG jest większy niż zbiór meczów z kursami, główne wyniki pracy raportowane
są na wspólnej próbce danych sportowych i rynkowych. Decyzja ta wynika z celu pracy, którym jest
porównanie modelu sportowego z rynkiem bukmacherskim. Tylko mecze posiadające kursy pozwalają
zestawić predykcje modelu, rankingów oraz benchmarku rynkowego na tych samych obserwacjach.

Pełny zbiór GOL.GG pozostaje jednak istotny, ponieważ dostarcza historycznej informacji potrzebnej
do budowania ratingów i cech kroczących. Ostateczna metodologia rozdziela zatem dwa zastosowania
danych: szeroka historia sportowa służy do konstruowania sygnałów predykcyjnych, natomiast wspólna
próbka danych sportowych i rynkowych służy do głównej ewaluacji i weryfikacji hipotez.
